\documentclass{article}

\usepackage{amsmath, amsthm, amssymb, amsfonts}
\usepackage{thmtools}
\usepackage{graphicx}
\usepackage{setspace}
\usepackage{geometry}
\usepackage{float}
\usepackage{hyperref}
\usepackage[utf8]{inputenc}
\usepackage[english]{babel}
\usepackage{framed}
\usepackage[dvipsnames]{xcolor}
\usepackage{tcolorbox}

\colorlet{LightGray}{White!90!Periwinkle}
\colorlet{LightOrange}{Orange!15}
\colorlet{LightGreen}{Green!15}

\newcommand{\HRule}[1]{\rule{\linewidth}{#1}}

\declaretheoremstyle[name=Theorem,]{thmsty}
\declaretheorem[style=thmsty,numberwithin=section]{theorem}
\tcolorboxenvironment{theorem}{colback=LightGray}

\declaretheoremstyle[name=Proposition,]{prosty}
\declaretheorem[style=prosty,numberlike=theorem]{proposition}
\tcolorboxenvironment{proposition}{colback=LightOrange}

\declaretheoremstyle[name=Principle,]{prcpsty}
\declaretheorem[style=prcpsty,numberlike=theorem]{principle}
\tcolorboxenvironment{principle}{colback=LightGreen}

\setstretch{1.2}
\geometry{
    textheight=9in,
    textwidth=5.5in,
    top=1in,
    headheight=12pt,
    headsep=25pt,
    footskip=30pt
}

% ------------------------------------------------------------------------------

\begin{document}

% ------------------------------------------------------------------------------
% Cover Page and ToC
% ------------------------------------------------------------------------------

\title{ \normalsize \textsc{}
		\\ [2.0cm]
		\HRule{1.5pt} \\
		\LARGE \textbf{\uppercase{Relatório do trabalho\\ Loja de Carros}
		\HRule{2.0pt} \\ [0.6cm] \LARGE{Ferramentas da Internet} \vspace*{10\baselineskip}}
		}
\date{}
\author{\textbf{Autores:} \\
		Mateus Augusto 1\\
		Kaique de Freitas 2\\
		Otavio Leone 3\\
        Lucas Estevam 4\\
        Tiago 5\\
        MARIANA - 2026}

\maketitle
\newpage

% ------------------------------------------------------------------------------

\section{Introdução}
O presente trabalho apresenta um diagrama entidade-relacionamento de um banco de dados para a loja de veículos. O objetivo principal do sistema é integrar, em uma estrutura única e coerente, os processos operacionais e comerciais do negócio: gestão de estoque, compras com fornecedores, vendas a clientes, financiamentos, agendamentos de test-drives e ordens de manutenção. Com as entidades:
Cliente,
Funcionário,
Venda,
Fornecedor,
Compra,
Financiamento,
Test-drive,
Manutenção,
Estoque,
Veiculo. 




\section{O que foi apresentado}
A documentação do projeto foi dividida em três pilares complementares que garantem a rastreabilidade dos requisitos:

Diagrama Entidade-Relacionamento (DER): Proporciona a visão espacial do banco de dados proporcionando um entendimento facilitado e dinamico.

Dicionário de Dados: É um documento ou repositório centralizado que descreve o significado, a estrutura, o tipo, a origem e o relacionamento dos elementos de dados em um sistema ou banco de dados.

Descrição de Relacionamentos: Define as regras rígidas de cardinalidade e dependência funcional entre as entidades.

\section{Conclusão}
A estrutura conceitual apresentada consolida de forma abrangente todas as regras de negócio extraídas do minimundo. O modelo garante a integridade referencial dos dados e fornece uma base sólida para a próxima etapa do projeto: o mapeamento para o modelo relacional (tabelas, chaves estrangeiras e normalização).

\newpage


\end{document}